% GENERATED FILE. Edit canonical lessons, then run npm run generate:books.
% canonical-source-hash: fnv1a64:708b62df4eced37d
% canonical-lessons: RU-C61-bird, RU-C61-mouse, RU-C61-wolf, RU-C61-bear, RU-C61-fox

\chapter{Bird, Mouse, Wolf, Bear, Fox}
\label{ch:ru-bird-mouse-wolf-bear-fox}
\hlchaptermodality{61}

\begin{chapteropening}
\textbf{By the end of this chapter:} \emph{I can name a bird, a mouse, a wolf, a bear and a fox.}

Complete the last lesson of chapter 61: I can name a bird, a mouse, a wolf, a bear and a fox.
\end{chapteropening}

\section[\texorpdfstring{\ru{птица} (ptítsa)}{птица (ptítsa)}]{\ru{птица} (ptítsa) — a bird}

\label{lesson:RU-C61-bird}



\begin{quote}
Before the new one: say the Russian for a cow, then the Russian for a pig.
\end{quote}

\subsection*{You'll want to know: \ru{птица}}
\textbf{\ru{птица}} (\emph{ptítsa}) — \textquotedblleft{}a bird\textquotedblright{}.

\begin{cousinweb}
\textbf{\ru{птица}} $\leftarrow$ Proto-Slavic \emph{*pŭtica}, from \emph{*pŭta}, \textquotedblleft{}bird\textquotedblright{} — inherited. Vasmer lists Latvian \emph{putns}, \textquotedblleft{}bird\textquotedblright{}, Lithuanian \emph{putýtis}, \textquotedblleft{}little bird\textquotedblright{}, and words for \textquotedblleft{}child, young one\textquotedblright{}: Latin \emph{putus}, \textquotedblleft{}boy\textquotedblright{}, and Sanskrit \emph{putrá-}, \textquotedblleft{}son\textquotedblright{}.
\end{cousinweb}

\begin{grammarlens}[title={What you've built}]
One more word for this chapter.
\end{grammarlens}

\subsection*{Your turn}
\begin{itemize}
\raggedright
  \item \emph{Say it:} \emph{ptítsa}
  \item \emph{Say it:} \emph{ptítsa}, once more
  \item \emph{Say it:} say \emph{ptítsa}
  \item \emph{Recall:} say the Russian for a cow, then the Russian for a pig, then say \emph{ptítsa} again
\end{itemize}

\begin{tcolorbox}[breakable,colback=teal!4,colframe=teal!35!black,title={Before you move on}]
What does \ru{птица} mean? (\textquotedblleft{}A bird\textquotedblright{}.) Say it once more.
\end{tcolorbox}
\section[\texorpdfstring{\ru{мышь} (mysh)}{мышь (mysh)}]{\ru{мышь} (mysh) — a mouse}

\label{lesson:RU-C61-mouse}



\begin{quote}
Before the new one: say the Russian for a pig, then the Russian for a bird.
\end{quote}

\subsection*{You'll want to know: \ru{мышь}}
\textbf{\ru{мышь}} (\emph{mysh}) — \textquotedblleft{}a mouse\textquotedblright{}.

\begin{cousinweb}
\textbf{\ru{мышь}} $\leftarrow$ Proto-Slavic \emph{*myšĭ} — inherited from PIE \emph{*muHs}, \textquotedblleft{}mouse\textquotedblright{}. Vasmer lists Sanskrit \emph{mūṣ-}, Greek \emph{mŷs}, Latin \emph{mūs} and Old High German \emph{mûs}: English \emph{mouse} is a cousin. Trubachev adds that it may first have meant \textquotedblleft{}the grey one\textquotedblright{}.
\end{cousinweb}

\begin{grammarlens}[title={What you've built}]
One more word for this chapter.
\end{grammarlens}

\subsection*{Your turn}
\begin{itemize}
\raggedright
  \item \emph{Say it:} \emph{mysh}
  \item \emph{Say it:} \emph{mysh}, once more
  \item \emph{Say it:} say \emph{mysh}
  \item \emph{Recall:} say the Russian for a pig, then the Russian for a bird, then say \emph{mysh} again
\end{itemize}

\begin{tcolorbox}[breakable,colback=teal!4,colframe=teal!35!black,title={Before you move on}]
What does \ru{мышь} mean? (\textquotedblleft{}A mouse\textquotedblright{}.) Say it once more.
\end{tcolorbox}
\section[\texorpdfstring{\ru{волк} (volk)}{волк (volk)}]{\ru{волк} (volk) — a wolf}

\label{lesson:RU-C61-wolf}



\begin{quote}
Before the new one: say the Russian for a bird, then the Russian for a mouse.
\end{quote}

\subsection*{You'll want to know: \ru{волк}}
\textbf{\ru{волк}} (\emph{volk}) — \textquotedblleft{}a wolf\textquotedblright{}.

\begin{cousinweb}
\textbf{\ru{волк}} $\leftarrow$ Proto-Slavic \emph{*vĭlkŭ} — inherited from PIE \emph{*wḷk\textsuperscript{w}os}, \textquotedblleft{}wolf\textquotedblright{}. Vasmer lists Lithuanian \emph{vilkas}, Sanskrit \emph{vṛkas}, Gothic \emph{wulfs} and Greek \emph{lýkos}: English \emph{wolf} is a cousin.
\end{cousinweb}

\begin{grammarlens}[title={What you've built}]
One more word for this chapter.
\end{grammarlens}

\subsection*{Your turn}
\begin{itemize}
\raggedright
  \item \emph{Say it:} \emph{volk}
  \item \emph{Say it:} \emph{volk}, once more
  \item \emph{Say it:} say \emph{volk}
  \item \emph{Recall:} say the Russian for a bird, then the Russian for a mouse, then say \emph{volk} again
\end{itemize}

\begin{tcolorbox}[breakable,colback=teal!4,colframe=teal!35!black,title={Before you move on}]
What does \ru{волк} mean? (\textquotedblleft{}A wolf\textquotedblright{}.) Say it once more.
\end{tcolorbox}
\section[\texorpdfstring{\ru{медведь} (medvéd)}{медведь (medvéd)}]{\ru{медведь} (medvéd) — a bear}

\label{lesson:RU-C61-bear}



\begin{quote}
Before the new one: say the Russian for a mouse, then the Russian for a wolf.
\end{quote}

\subsection*{You'll want to know: \ru{медведь}}
\textbf{\ru{медведь}} (\emph{medvéd}) — \textquotedblleft{}a bear\textquotedblright{}. Literally \textquotedblleft{}honey-knower\textquotedblright{}: the one who knows where the \textbf{\ru{мёд}} is.

\begin{cousinweb}
\textbf{\ru{медведь}} $\leftarrow$ Proto-Slavic \emph{*medvědĭ}, \textquotedblleft{}honey-eater\textquotedblright{}, from \emph{\ru{мёд}}, \textquotedblleft{}honey\textquotedblright{}, and the root of \emph{\ru{есть}}, \textquotedblleft{}to eat\textquotedblright{} (Vasmer). It is a taboo name: hunters avoided the old word for \textquotedblleft{}bear\textquotedblright{}, the one kept in Latin \emph{ursus} and Greek \emph{árktos}.
\end{cousinweb}

\begin{grammarlens}[title={What you've built}]
One more word for this chapter.
\end{grammarlens}

\subsection*{Your turn}
\begin{itemize}
\raggedright
  \item \emph{Say it:} \emph{medvéd}
  \item \emph{Say it:} \emph{medvéd}, once more
  \item \emph{Say it:} say \emph{medvéd}
  \item \emph{Recall:} say the Russian for a mouse, then the Russian for a wolf, then say \emph{medvéd} again
\end{itemize}

\begin{tcolorbox}[breakable,colback=teal!4,colframe=teal!35!black,title={Before you move on}]
What does \ru{медведь} mean? (\textquotedblleft{}A bear\textquotedblright{}.) Say it once more.
\end{tcolorbox}
\section[\texorpdfstring{\ru{лиса} (lisá)}{лиса (lisá)}]{\ru{лиса} (lisá) — a fox}

\label{lesson:RU-C61-fox}



\begin{quote}
Before the new one: say the Russian for a wolf, then the Russian for a bear.
\end{quote}

\subsection*{You'll want to know: \ru{лиса}}
\textbf{\ru{лиса}} (\emph{lisá}) — \textquotedblleft{}a fox\textquotedblright{}.

\begin{cousinweb}
\textbf{\ru{лиса}} goes with \emph{\ru{лис}}, \textquotedblleft{}a fox\textquotedblright{}, Old Church Slavonic \emph{\ru{лиса}}. Vasmer calls the origin unclear: tracing it to the family of Latin \emph{vulpēs} and Lithuanian \emph{lãpė}, \textquotedblleft{}fox\textquotedblright{}, is guesswork, he says, and other links are doubtful.
\end{cousinweb}

\begin{grammarlens}[title={What you've built}]
One more word for this chapter.
\end{grammarlens}

\subsection*{Your turn}
\begin{itemize}
\raggedright
  \item \emph{Say it:} \emph{lisá}
  \item \emph{Say it:} \emph{lisá}, once more
  \item \emph{Say it:} say \emph{lisá}
  \item \emph{Recall:} say the Russian for a wolf, then the Russian for a bear, then say \emph{lisá} again
\end{itemize}

\begin{tcolorbox}[breakable,colback=teal!4,colframe=teal!35!black,title={Before you move on}]
What does \ru{лиса} mean? (\textquotedblleft{}A fox\textquotedblright{}.) Say it once more.
\end{tcolorbox}
